
\documentclass{article}
\usepackage{colortbl}
\usepackage{makecell}
\usepackage{multirow}
\usepackage{supertabular}

\begin{document}

\newcounter{utterance}

\centering \large Interaction Transcript for game `wordle', experiment `medium\_frequency\_words\_no\_clue\_no\_critic', episode 2 with reflexion{-}Qwen3.5{-}9B.
\vspace{24pt}

{ \footnotesize  \setcounter{utterance}{1}
\setlength{\tabcolsep}{0pt}
\begin{supertabular}{c@{$\;$}|p{.15\linewidth}@{}p{.15\linewidth}p{.15\linewidth}p{.15\linewidth}p{.15\linewidth}p{.15\linewidth}}
    \# & \multicolumn{2}{c}{Player} && \multicolumn{2}{c}{Game Master} \\
    \hline

    \theutterance \stepcounter{utterance}  
    & & & \multicolumn{4}{p{0.6\linewidth}}{
        \cellcolor[rgb]{0.9,0.9,0.9}{
            \makecell[{{p{\linewidth}}}]{
                \texttt{\tiny{[P1$\langle$GM]}}
                \texttt{You have attempted this same task before and did not succeed. These are the plans you wrote after those attempts.} \\
\texttt{Trial \#1: New plan: I failed because I misinterpreted the feedback from "renew". The feedback `n<yellow>` meant 'n' is in the word but NOT at index 2. I incorrectly assumed 'n' was at index 2 in "renew" (it is actually at index 2 in the word "renew", so the yellow correctly indicated it's in the wrong spot relative to the target, but I got confused about the target structure). More critically, I wasted moves guessing "nerve", "nerdy", and "renew" without systematically filling the remaining two slots (indices 2 and 3) using high{-}frequency consonants like 's', 't', 'l', 'd' while keeping 'r', 'e', 'n' fixed in their known valid positions (r at 0, e at 1, n somewhere in 3 or 4). Next time, once I have `r<green> e<green>`, I will immediately test "rents" or "reins" to place 'n' at the end or 's'/'t' in the middle, rather than guessing words that repeat 'e' or 'n' in positions I already know are wrong. I will strictly track the set of confirmed letters and their forbidden indices to generate the next guess.} \\
\\ 
\\ 
\texttt{You are a language wizard who likes to guess words by using the given rules.} \\
\\ 
\texttt{Welcome to Wordle! You have six attempts to guess the target word, a valid English word of five lowercase letters (a{-}z). Please use the tags "explanation:" and "guess:" to provide a concise explanation for each guess.} \\
\\ 
\texttt{For instance, if your guess is "apple", your response should be} \\
\texttt{explanation: this is a common five{-}letter English word, and I am starting my guess with this word.} \\
\texttt{guess: apple} \\
\\ 
\texttt{After each guess, your answer will be validated, and you will receive feedback indicating which letters are correct (green), which letters are correct but in the wrong position (yellow), and which letters are incorrect (red). This feedback can be useful in determining which letters to include or exclude in your next guess.} \\
\\ 
\texttt{For example, the feedback for "apple" might be:} \\
\texttt{guess\_feedback: a<yellow> p<yellow> p<green> l<yellow> e<red>} \\
\\ 
\texttt{The explanation should contain details about how the guess\_feedback is used to arrive at a new guess.} \\
\\ 
\texttt{Do not generate any other text.} \\
\\ 
\texttt{Let's begin with your first guess.} \\
            }
        }
    }
     \\ \\

    \theutterance \stepcounter{utterance}  
    & \multicolumn{4}{p{0.6\linewidth}}{
        \cellcolor[rgb]{0.561,0.737,0.561}{
            \makecell[{{p{\linewidth}}}]{
                \texttt{\tiny{[P1$\rangle$GM]}}
                \texttt{explanation: I am starting with "crane" as my first guess because it contains common vowels (a, e) and high{-}frequency consonants (r, n) that appear in many five{-}letter words. This will help me quickly identify which letters are present and their potential positions.} \\
\texttt{guess: crane} \\
            }
        }
    }
    & & \\ \\

    \theutterance \stepcounter{utterance}  
    & & & \multicolumn{2}{p{0.3\linewidth}}{
        \cellcolor[rgb]{0.9,0.9,0.9}{
            \makecell[{{p{\linewidth}}}]{
                \texttt{\tiny{[GM$|$GM]}}
                \texttt{attempts = 1} \\
\texttt{target\_word = resin} \\
\texttt{guess = crane} \\
\texttt{guess\_feedback = c<red> r<yellow> a<red> n<yellow> e<yellow>} \\
            }
        }
    }
    & & \\ \\

    \theutterance \stepcounter{utterance}  
    & & & \multicolumn{4}{p{0.6\linewidth}}{
        \cellcolor[rgb]{0.9,0.9,0.9}{
            \makecell[{{p{\linewidth}}}]{
                \texttt{\tiny{[P1$\langle$GM]}}
                \texttt{guess\_feedback: c<red> r<yellow> a<red> n<yellow> e<yellow>} \\
\\ 
\texttt{Provide your response only in this format.} \\
\texttt{explanation: details} \\
\texttt{guess: word} \\
            }
        }
    }
     \\ \\

    \theutterance \stepcounter{utterance}  
    & \multicolumn{4}{p{0.6\linewidth}}{
        \cellcolor[rgb]{0.561,0.737,0.561}{
            \makecell[{{p{\linewidth}}}]{
                \texttt{\tiny{[P1$\rangle$GM]}}
                \texttt{explanation: The feedback indicates that 'r', 'n', and 'e' are in the word but in wrong positions. Specifically, 'r' is not at index 1, 'n' is not at index 3, and 'e' is not at index 4. 'c' and 'a' are not in the word. I will keep 'r', 'n', and 'e' as candidates and try to place them in different positions while introducing high{-}frequency letters like 's', 't', or 'l' to fill the remaining slots. I'll guess "rents" to test if 'n' belongs at index 3 (which it shouldn't based on previous feedback, so maybe 'n' is elsewhere) and to check 's' and 't'. Wait, 'n' was yellow at index 3 in "crane", meaning 'n' is in the word but NOT at index 3. So "rents" puts 'n' at index 1, which is fine. Let's try "rents" to see if 'e' moves or if 's'/'t' are correct.} \\
\texttt{guess: rents} \\
            }
        }
    }
    & & \\ \\

    \theutterance \stepcounter{utterance}  
    & & & \multicolumn{2}{p{0.3\linewidth}}{
        \cellcolor[rgb]{0.9,0.9,0.9}{
            \makecell[{{p{\linewidth}}}]{
                \texttt{\tiny{[GM$|$GM]}}
                \texttt{attempts = 2} \\
\texttt{target\_word = resin} \\
\texttt{guess = rents} \\
\texttt{guess\_feedback = r<green> e<green> n<yellow> t<red> s<yellow>} \\
            }
        }
    }
    & & \\ \\

    \theutterance \stepcounter{utterance}  
    & & & \multicolumn{4}{p{0.6\linewidth}}{
        \cellcolor[rgb]{0.9,0.9,0.9}{
            \makecell[{{p{\linewidth}}}]{
                \texttt{\tiny{[P1$\langle$GM]}}
                \texttt{guess\_feedback: r<green> e<green> n<yellow> t<red> s<yellow>} \\
\\ 
\texttt{Provide your response only in this format.} \\
\texttt{explanation: details} \\
\texttt{guess: word} \\
            }
        }
    }
     \\ \\

    \theutterance \stepcounter{utterance}  
    & \multicolumn{4}{p{0.6\linewidth}}{
        \cellcolor[rgb]{0.561,0.737,0.561}{
            \makecell[{{p{\linewidth}}}]{
                \texttt{\tiny{[P1$\rangle$GM]}}
                \texttt{explanation: The feedback confirms 'r' is at index 0 and 'e' is at index 1. 'n' is in the word but not at index 2 (it was yellow in "rents" at index 2). 's' is in the word but not at index 4. 't' is not in the word. With "r\_e\_n\_s" pattern where 'n' cannot be at index 2 and 's' cannot be at index 4, I need to place 'n' at index 3 or 4. Since 's' is yellow at index 4, 's' might be at index 3. If 'n' is at index 3, then 's' must be elsewhere, but 's' is already confirmed. Let's try "renos" to test if 'o' is the word or if 'n' is at index 3 and 's' is at index 4 (which is forbidden). Wait, 's' is yellow at index 4 in "rents", so 's' is not at index 4. In "renos", 's' is at index 4 again, which is bad. Let's try "rings" to test 'i' and 'g' while keeping 'r', 'e', 'n', 's' in positions 0,1,3,4? No, 'n' was yellow at index 2 in "rents", so 'n' is not at index 2. In "rings", 'n' is at index 2. Bad. Let's try "rises" {-}> r(0), i(1), s(2), e(3), s(4). 'e' must be at 1. Bad. Let's try "rines" {-}> r(0), i(1), n(2), e(3), s(4). 'n' at 2 is bad. Let's try "rones" {-}> r(0), o(1), n(2), e(3), s(4). 'e' at 3 is bad. The only valid positions for 'n' are 3 or 4. If 'n' is at 3, the word is r\_e\_n\_s. If 'n' is at 4, the word is r\_e\_s\_n. But 's' is yellow at 4 in "rents", so 's' is not at 4. If the word is r\_e\_s\_n, 's' is at 2. Let's try "resin" to test 'i' and 's' at 2, 'n' at 4.} \\
\texttt{guess: resin} \\
            }
        }
    }
    & & \\ \\

    \theutterance \stepcounter{utterance}  
    & & & \multicolumn{2}{p{0.3\linewidth}}{
        \cellcolor[rgb]{0.9,0.9,0.9}{
            \makecell[{{p{\linewidth}}}]{
                \texttt{\tiny{[GM$|$GM]}}
                \texttt{attempts = 3} \\
\texttt{target\_word = resin} \\
\texttt{guess = resin} \\
\texttt{guess\_feedback = r<green> e<green> s<green> i<green> n<green>} \\
            }
        }
    }
    & & \\ \\

    \theutterance \stepcounter{utterance}  
    & & & \multicolumn{2}{p{0.3\linewidth}}{
        \cellcolor[rgb]{0.9,0.9,0.9}{
            \makecell[{{p{\linewidth}}}]{
                \texttt{\tiny{[GM$|$GM]}}
                \texttt{game\_result = WIN} \\
            }
        }
    }
    & & \\ \\

\end{supertabular}
}

\end{document}
